\documentclass[twocolumn,a4paper,10pt]{article}
\usepackage[utf8]{inputenc}
\usepackage[english]{babel}
\usepackage{amsmath}
\usepackage{graphicx}
\usepackage{geometry}
\geometry{a4paper, left=2cm, right=2cm, top=2.5cm, bottom=2.5cm}
\usepackage{authblk}

\title{Analysis of the Evolution of the Upper Tropospheric Cyclonic Vortex of December 1995 over Southeastern South America based on Isentropic Potential Vorticity}
\author{Author Name}
\affil{Institution}
\date{}

\begin{document}

\twocolumn[
  \begin{@twocolumnfalse}
    \maketitle
    \begin{abstract}
During the last two weeks of 1995, a baroclinic system crossed the Andes and moved slowly northeastward over the coast of South and Southeast of Brazil as an Upper Tropospheric Cyclonic Vortex. It produced an enormous amount of precipitation, flash floods and great human and material damages. In this work, the anomalous path and the process of deepening and decaying of the system is investigated based on the isentropic potential vorticity - IPV (IPV thinking). It is shown that the baroclinic instability played an important role in its formation west of the Andes. After crossing the Andes, it induced a lee cyclogenesis. The interaction with a surface temperature disturbance in southern Brazil intensified the IPV anomaly, but not accompanied by a new baroclinic instability development. The anomalous path to northeast occurs with strong north-south shear of zonal wind where the necessary condition for the barotropic instability is satisfied.
\vspace{0.5cm} \\
\textbf{Key words:} Upper Tropospheric Cyclonic Vortex; IPV thinking; Flash Flood.
\vspace{1cm}
    \end{abstract}
  \end{@twocolumnfalse}
]

\section{INTRODUCTION}
The study of the atmospheric dynamics responsible for the development of extratropical cyclones -- hereafter EC, is of notorious interest to meteorology (Bjerknes H.B, 1922; Charney, J. G., 1947). This is due both to the role that ECs play in atmospheric energetics and to the changes in weather conditions associated with them.

Equally important, the concept of Potential Vorticity-PV (Ertel H, 1942 and Rossby 1937b), has been widely used in the study of ECs since the works of Kleinschmidt (1957). The introduction of the concept of Isentropic Potential Vorticity (IPV) by Hoskins, B.J, et al. (1985), gave a new impetus to the use of this technique. According to Thorpe, A.J. (1985), due to this new way of thinking about the atmosphere, this technique was called "IPV Thinking" (hereafter referred to as IPV analysis).

In this article, IPV analysis is used to study the dynamics that produced an Extratropical Cyclone - EC, over Southeastern South America, between December 10 and 31, 1995. Its long duration, the anomalous trajectory to the northeast, and several heavy rain events on the coast of Southern and Southeastern Brazil justify an investigation of this phenomenon.

\section{LITERATURE REVIEW}
Baroclinic instability is the main mechanism for Extratropical Cyclone-EC formation (Hakim, G. J.L. et al, 1995). In this context, the surface pressure drop is mainly due to the warming caused by the extrusion of stratospheric air in the upper and middle troposphere (Hirshberg, P.A., and M.J. Fritsch, 1991a and 1995b). Thus, in an environment with vertical wind shear, the extrusion of stratospheric air forms a tropopause fold (or an IPV tower) which amplifies downwards in an atmosphere with strong vertical wind shear.

In South America, the trajectory of EC displacement has been studied by Necco G. V., 1982; Satyamurty, P., Ferreira, C.C., Gan, M.A, (1990); Sinclair M.R. (1995); Ferreira, C.C., (1989) and Ambrizzi, T., and A. B. Pezza, (1999). In general, the authors show that ECs present a preferential trajectory to the east and southeast.

The role of the Andes relief in the development of ECs was studied by Orlanski, I., J., et al. (1991). They carried out sensitivity tests with a limited area model of an EC event over southern South America. The results showed that its development was insensitive to the presence of the Andes relief and surface fluxes. However, surface friction substantially affected the EC intensification.

The role of the Andes Cordillera in the development of cyclogenesis on the coast of Uruguay in July 1996 was studied by Funatsu, B. M., et al. (2004). The authors classified the event as lee cyclogenesis, since the existence of the mountain barrier was fundamental for the intensification of surface baroclinity, which, in turn, led to the formation of the extratropical cyclone.

Reboita, M. S. et al. (2009) analyzed the quasi-stationarity of a Cyclone on the Southern Coast of Brazil at the beginning of May 2008, using Sutcliffe's equation. According to the author, the important mechanism for the system's cyclogenesis was the advection of cyclonic absolute vorticity in the middle troposphere and warm air advection in the 1000-500 hPa layer. This advection was associated with an Upper Tropospheric Cyclonic Vortex (UTCV) that formed in a region of potential vorticity anomaly. The UTCV remained semi-stationary and composed the northern sector of a dipole-type blocking. The atypical movement of the UTCV was associated with advections of cyclonic absolute vorticity in the middle troposphere and warm air in its southern sector. Only when the blocking at middle levels and the potential vorticity anomaly at middle/upper levels weakened, did the surface cyclone move away from the southern coast of Brazil.

Iwabe, C. M. et al. (2010) carried out a case study of the atmospheric situation of an extratropical cyclone between September 12 and 19, 2008, and the Catarina event. When moving eastwards, the main similarities between the event under study and Hurricane Catarina were the presence of a dipole-type blocking pattern and upper-level PV anomalies with a westward tilt.

The blocking associated with the September 2008 event lasted a day and a half, while the Hurricane Catarina blocking lasted three days. According to the authors, in the September 2008 event, the warm advection to the east of the system seems to have contributed to shifting it eastwards, while in the Catarina event, a drastic decrease in the warm advection to the east of the system was observed.

\sectionsection{MATERIALS AND METHODS}
NCEP (National Center for Environmental Prediction) Reanalysis I data (Kalnay et al., 1996) were used for the period from December 10 to 31, 1995. The data have a horizontal grid spacing of 2.5$^{\circ}$ latitude by 2.5$^{\circ}$ longitude, with 144x73 global grid points, at 17 vertical pressure levels (1000, 925, 850, 700, 600, 500, 400, 300, 250, 200, 150, 100, 70, 50, 30, 20, 10 hPa) and at four daily times (00:00, 06:00, 12:00 and 18:00 UTC).

The dynamics analysis is performed using isentropic potential vorticity -- IPV (Hoskins, et al., 1985):

\begin{equation}
IPV = g (\zeta_{\theta} + f) \left( - \frac{\partial \theta}{\partial p} \right)
\end{equation}

where $g$ is the acceleration of gravity, $f = 2\Omega \sin(\phi)$ is the Coriolis parameter; $\Omega = 7.292 \times 10^{-5}$ rad/s is the earth's angular velocity and $\phi$ is the local latitude; $S = - \frac{\partial \theta}{\partial p}$ is the static stability; $\zeta_{\theta}$ is the relative vorticity in isentropic coordinates, $\theta$ is the potential temperature and $p$ is the air pressure.

A potential vorticity unit (PVU=$10^{-6}$ m$^{2}$ K kg$^{-1}$ s$^{-1}$) should be associated with a 10K change in potential temperature for a vertical displacement of 100 hPa (Bluestein, H.B., 1993). In this article, the cyclonic and anti-cyclonic IPV anomaly convention will be adopted, as this convention is valid in both hemispheres.

IPV is conserved in an isentropic flow. The conservation of IPV implies that the static stability and absolute vorticity mutually compensate each other.

According to (Bluestein, H.B., 1993), the stratosphere is an infinite reservoir of cyclonic IPV, so the main mechanism of cyclonic IPV production in the upper and middle troposphere is the vertical and horizontal advection of IPV. In the lower troposphere, IPV is created by differential diabatic heating and interaction with mountains. IPV can only be destroyed by diabatic processes (latent and sensible heat) and by friction. Friction acts on the wind, causing momentum to not be conserved. The release of latent heat, associated with precipitation, warms the air parcel and decreases static stability. When it falls, the precipitation "transports" the cyclonic IPV towards the surface.

A conceptual model of EC development was suggested by Hirshberg, P.A., and M.J. Fritsch (1991b). According to this model, in a baroclinic atmosphere, the presence of a cyclonic IPV anomaly at higher levels induces a cyclonic IPV anomaly at the surface, in a region at L/4 of the wavelength L, associated with the upper anomaly, and downstream of its position. The same amplification should occur if the system phases with a surface temperature disturbance (seen in terms of the equivalent potential temperature increase).

Hirshberg, P.A., and M.J. Fritsch (1991b) show that if a low-pressure core exists at the surface, it responds to the warming of air at upper and middle levels. Mathematically, this is shown by taking the surface pressure tendency equation:

\begin{equation}
\frac{\partial (p_s)}{\partial t} = \frac{R}{g} \int \frac{\partial T_v}{\partial t} d(\ln p)
\end{equation}

Where $p_s$ represents the surface pressure, $p$ the pressure, $T_v$ the virtual temperature, $R$ is the gas constant and $g$ is the acceleration of gravity. The surface pressure tendency depends on the heating, or $\frac{\partial T_v}{\partial t}$, integrated vertically. Note that this term is multiplied by $d(\ln p)$. Thus, heating at upper levels is more effective than near the surface.

To verify the model suggested by Hirshberg, P.A., and M.J. Fritsch (1991b), vertical profiles of IPV and potential temperature (isentropes) are plotted. In these profiles, the warmer air regions are perceived by the downward inclination of the isentropes and the opposite for colder regions.

The investigation of the northeastward trajectory is done from the perspective of the development of barotropic instability. The necessary condition for barotropic instability is obtained when the term:

\begin{equation}
Q_y = \beta - \frac{d^2 \bar{u}}{dy^2}
\end{equation}

changes sign within a certain domain (Kuo, 1949), where $\beta = \frac{df}{dy}$ is the variation of the Coriolis parameter with distance $y$, in the north-south direction and $\bar{u}$ is the mean zonal wind.

\section{RESULTS}
The EC event under study in this article developed between December 10 and 31, 1995. In this period, the upper-level perturbation underwent several transformations associated with the presence of IPV anomalies at the surface, the presence of the Andes Cordillera and the presence of a surface temperature disturbance. To facilitate identification, the stages of the system under study can be divided into:

\begin{enumerate}
    \item UTCV/EC Formation Phase (10 to 14/12/1995);
    \item Rapid Intensification Phase, associated with baroclinic instability (14 to 19/12/1995);
    \item Phase associated with the crossing of the Andes Cordillera, where a strong lee cyclogenesis was observed in northern Argentina (21/12/1995);
    \item Mature Phase; where a new intensification of the surface cyclone was observed (22 to 23/12/1995);
    \item Decay Phase (27 to 30/12/1995).
\end{enumerate}

\subsection{UTCV Formation Stage}
The formation process of the UTCV/EC occurred over the Pacific Ocean between 10 and 14/12/1995. This vortex originated from a cold front (not shown) that had its equatorial part moving more slowly than the polar one. This type of evolution agrees with the traditional view of UTCV formation (Gan, M. A., and V. B. Rao, 1994). In the GOES-8 satellite image, an already formed vortex with little associated cloudiness can be verified on 14/12/1995 and 19/12/1995 at 17:45 UTC.

\subsection{Baroclinic Intensification Stage}
From 18/12/1995, this upper-level cyclonic IPV anomaly showed rapid intensification in a typically baroclinic process. This process occurred between 18 and 20/12/1995, when this upper-level perturbation approached a small surface cyclonic IPV anomaly, located around 80W and 35S. This latter surface cyclonic IPV anomaly was dammed in this area by the Andes Cordillera. The rapid intensification of the cyclonic IPV anomalies at upper and lower levels occurred when there was an overlap of the two cyclonic anomalies, in an L/4 relationship, where L is the wavelength. The cyclonic vorticity anomaly near the surface was dammed west of the Andes since the passage of a previous system. This anomaly amplified when it came into phase with the upper-level anomaly. This in-phase amplification agrees with the conceptual model of baroclinic amplification by Hirshberg, P.A., and M.J. Fritsch (1991b) and is due to differential temperature advection (Hoskins et al 1995).

\subsection{Lee Cyclogenesis of the Andes}
A characteristic frequently observed in studies of UTCVs that cross mountains like the Andes is a change in the vertical inclination of the system (Gan, M. A.; Rao, V. B. 1994; Funatsu, B. M., et al., 2004). It is possible to verify that its inclination changed from west to east and, finally, it assumed an equivalent barotropic vertical structure.
At the end of 20/12/1995 and early morning of 21/12/1995, a lee cyclogenesis occurred east of the Andes. According to data in Mendoza, the pressure had dropped 18 hPa in 24 hours and was at 994 hPa at 00Z 21/12/1995. During the cyclone formation, strong winds were observed over the Argentine Pampa.
Two relatively warmer regions can be verified over the low-pressure center. The first warm region is near the surface, east of the Andes. The second warm region is associated with the UTCV between 100 and 300 hPa. According to Hirschberg and Fritsch (1991b), the warming at high levels is a result of stratospheric air extrusion.
In the early morning following the lee cyclogenesis and in the same area, a strong cooling of the air is seen. Along with cooling, it was observed that the surface low-pressure center rapidly lost strength (1005 hPa).

\subsection{Mature Phase of the System}
After crossing the Andes, the IPV tower presented an equivalent barotropic structure. A surface temperature disturbance is visible over Uruguay and the South and Center of the state of Rio Grande do Sul. From the point of view of IPV analysis, this surface temperature disturbance should act to generate a cyclonic IPV anomaly near the surface (Hoskins, B.J, et al, 1985). Concomitantly, an amplification of the IPV tower occurred and a new deepening of the surface low-pressure center went from 1005 to 1000 hPa. 
This stage was accompanied by heavy precipitation over the eastern sector of the UTCV. On the coast of the states of RS, SC, PR, and SP, the heaviest precipitation was recorded. The disappearance of the cyclonic IPV anomaly in the region with the most precipitation is compatible with the arguments of (Hoskins, B.J, et al, 1985), regarding the destruction of cyclonic IPV by rain.

\subsection{Decay Stage}
The beginning of the decay stage occurred at 12 UTC on 28/12/1995. From this phase, the surface pressure center was no longer visible. Heavy precipitation events were still recorded on the coast of Paraná, São Paulo, and Rio de Janeiro.
In the IPV fields, it was possible to track this UTCV until 12 UTC on January 1, 1996. After this day, the UTCV was finally incorporated into a cold front that was forming over the Southeast.

\subsection{Northeastward Trajectory}
After 22/12/1995, the UTCV began to move slowly northeastward. Normally, UTCVs and ECs move southeast and east (Ambrizzi, T., and A. B. Pezza, 1999). A possible explanation for this anomalous displacement may stem from barotropic instability. From 29/12/1995, the north-south shear of the zonal wind rapidly disappears and with it $Q_y$. According to (Kuo, A. L., 1949), the necessary condition for barotropic instability is the change of sign of this parameter $Q_y$.
The north-south shear of the zonal wind is necessary for the existence of barotropic instability. Apparently, the formation of this north-south shear was linked to the existence of the cold low (UTCV); however, the existence of a high-pressure system over the Ocean, near the surface, may have contributed.

\section{FINAL CONSIDERATIONS}
In this article, a case study was made of a UTCV, of long duration and anomalous northeastward trajectory, which caused floods during the Christmas period of 1995 over the Southern and Southeastern Coast of Brazil. The technique used was IPV analysis, or "IPV Thinking" by Thorpe, A.J., (1985) and a conceptual model of cyclone development suggested by Hirshberg, P.A., and M.J. Fritsch (1991b).
Comparing the conceptual model with the IPV vertical profiles, it was found that the system's intensification west of the Andes was due to baroclinic amplification, in an L/4 phase coupling of an upper-level cyclonic IPV anomaly with a near-surface cyclonic anomaly.
The strong lee cyclogenesis over Northern Argentina on 21/12/1995 was due to the overlap of the warmer air column, associated with the UTCV, with the warm air existing over the Continent. 
The surface low-pressure center always accompanied the center of the IPV tower at 400 hPa. Since 23/12/1995, the IPV tower started to present an inverse vertical inclination, towards the east. The decay of the IPV tower was mainly due to precipitation. 
A possible explanation for the system's anomalous northeastward trajectory can be made based on barotropic instability. The change in sign of the parameter $Q_y$ is a necessary condition for the existence of barotropic instability. Another important aspect to explain the long duration of the UTCV, as well as its northeastward trajectory, is the existence of a dipole block.

\section{REFERENCES}
\begin{itemize}
    \item Ambrizzi, T., and A. B. Pezza, 1999: Cold waves and the propagation of extratropical cyclones and anticyclones in South America. Revista Geofisica, 51, 45-67.
    \item Bjerknes, J., and Solberg, H., 1922: Life cycles of cyclones and the polar front theory of atmospheric circulation, Geofys. Publ., 3(1), 1-18.
    \item Bluestein, H.B., 1993: Synoptic-dynamic meteorology in the midlatitudes. Volume II: Observations and theory of weather systems. Oxford University Press, New York, 594pp.
    \item Charney, J. G. (1947). The dynamics of long waves in a baroclinic westerly current. J. Meteor., 4:135-163.
    \item Funatsu, B. M., et al, 2004: A case study of orographic cyclogenesis over South America. Atmósfera, 91-113.
    \item Gan, M. A., and V. B. Rao, 1994. The influence of the Andes Cordillera on transient disturbances. Mon. Wea. Rev., 122 (6), 1141-1157.
    \item Hirshberg, P.A., and M.J. Fritsch, 1991a: Tropopause undulations and the development of extratropical cyclones. Part I. Mon.Wea.Rev., 119, 496-517.
    \item Hirshberg, P.A., and M.J. Fritsch, 1991b: Tropopause undulations and the development of extratropical cyclones. Part II. Mon.Wea.Rev., 119, 518-550.
    \item Hoskins, B.J., M.E. McIntyre and A.W. Robertson, 1985: On the use and significance of isentropic potential vorticity maps. Quart.J.Roy.Meteor.Soc., 111, 877-946.
\end{itemize}

\end{document}
